% notes-metric-spaces-with-mathlib.tex
% Topic: Metric Spaces with Mathlib

\begin{remark*}[Exposition]
This topic records the mathlib-based formalization track for metric spaces.
The mathematical definition below is a non-extracting reminder: in Lean with
mathlib, the relevant structure is already supplied by the library, so the
formalization task is to learn how to use and align with that existing API.
\end{remark*}

\begin{remark*}[Mathematical reminder]
\begin{definition*}[Metric Space]
A \textbf{metric space} is a pair $(X,d)$ consisting of a set $X$ and a
function
\[
  d\colon X\times X \to [0,\infty)
\]
such that for all $x,y,z\in X$:
\begin{enumerate}
  \item $d(x,y)=0$ if and only if $x=y$;
  \item $d(x,y)=d(y,x)$;
  \item $d(x,z)\leq d(x,y)+d(y,z)$.
\end{enumerate}
\end{definition*}
\end{remark*}

\begin{remark*}[Formalization note]
The active Lean development is Mathlib-first. Main theorem files work in the
context
\[
  \texttt{variable \{X : Type u\} [MetricSpace X]}.
\]
The distance function is Mathlib's global notation \texttt{dist x y}, and the
basic metric facts are the existing API lemmas \texttt{dist\_self},
\texttt{dist\_comm}, \texttt{dist\_nonneg}, and \texttt{dist\_triangle}. The
old bundled metric structures have been renamed to \texttt{ScratchMetric} and
\texttt{ScratchMetricSpace}; they are retained only as pedagogical scratch
models for exposing raw proof obligations.
\end{remark*}

\subsection{Distance to Sets and Diameter}

\begin{remark*}[Mathematical reminder]
\begin{definition*}[Distance from a Point to a Set]
Let $(X,d)$ be a metric space, let $x\in X$, and let $A\subseteq X$ be
nonempty. The \textbf{distance from \(x\) to \(A\)} is
\[
  d(x,A)=\inf\{d(x,a)\mid a\in A\}.
\]
\end{definition*}
\end{remark*}

\begin{remark*}[Lean definition]
The Lean file \texttt{SetDistance.lean} names the set of candidate distances:
\begin{verbatim}
def distanceSet
    {X : Type u}
    [MetricSpace X]
    (x : X)
    (S : Set X) : Set Real :=
  (fun y : X => dist x y) '' S
\end{verbatim}
The point-to-set distance is then the real-valued infimum of that set:
\begin{verbatim}
noncomputable def distanceToSet
    {X : Type u}
    [MetricSpace X]
    (x : X)
    (S : Set X) : Real :=
  sInf (distanceSet x S)
\end{verbatim}
This is deliberately not a new structure. It is a function of the ambient
Mathlib metric-space instance, the point, and the set.
\end{remark*}

\begin{remark*}[Formalization issue]
This definition introduces the first order-theoretic object in the metric
track: \texttt{sInf}. The current development is real-valued, so the empty-set
case should not be read as the extended-real convention
\(d(x,\emptyset)=+\infty\). The proof work is therefore
organized around explicit hypotheses such as \texttt{S.Nonempty} and around
helper lemmas proving that \texttt{distanceSet x S} is nonempty and bounded
below.
\end{remark*}

\begin{remark*}[Mathematical reminder]
\begin{definition*}[Diameter of a Set]
Let $(X,d)$ be a metric space and let $A\subseteq X$ be nonempty. The
\textbf{diameter} of \(A\) is
\[
  \operatorname{diam}(A)
  =
  \sup\{d(a,b)\mid a,b\in A\}.
\]
\end{definition*}
\end{remark*}

\begin{remark*}[Lean definition]
The Lean file \texttt{SetDiameter.lean} similarly names the set whose
supremum is being taken:
\begin{verbatim}
def diameterSet
    {X : Type u}
    [MetricSpace X]
    (S : Set X) : Set Real :=
  { r : Real | exists x : X, x in S and
      exists y : X, y in S and r = dist x y }
\end{verbatim}
The diameter itself is
\begin{verbatim}
noncomputable def diameter
    {X : Type u}
    [MetricSpace X]
    (S : Set X) : Real :=
  sSup (diameterSet S)
\end{verbatim}
\end{remark*}

\begin{remark*}[Formalization issue]
The diameter definition brings in \texttt{sSup}. The current monotonicity
theorem is proved by separating the metric-specific and order-theoretic parts:
\texttt{diameterSet\_mono} turns \(A\subseteq B\) into inclusion of pairwise
distance sets, and \texttt{csSup\_le\_csSup} turns that inclusion into an
inequality between suprema when the required nonemptiness and bounded-above
hypotheses are supplied.
\end{remark*}

\begin{remark*}[Proof targets]
\begin{theorem*}[A Singleton Has Diameter Zero]
Let $(X,d)$ be a metric space and let \(x\in X\). Then
\[
  \operatorname{diam}(\{x\})=0.
\]
\end{theorem*}
\end{remark*}

\begin{remark*}[Proof target]
\phantomsection
\label{thm:fv-msml-diameter-monotone-under-inclusion}
\begin{theorem*}[Diameter Is Monotone]
Let $(X,d)$ be a metric space. If \(A\subseteq B\), the pairwise distance set
for \(A\) is nonempty, and the pairwise distance set for \(B\) is bounded
above, then
\[
  \operatorname{diam}(A)\leq \operatorname{diam}(B).
\]
\end{theorem*}
\hyperref[prf:fv-msml-diameter-monotone-under-inclusion]{\textit{Go to proof.}}
\end{remark*}

\begin{remark*}[Proof target]
\phantomsection
\label{thm:fv-msml-distance-to-set-le-distance-to-point-of-mem}
\begin{theorem*}[Point-to-Set Distance Is Bounded by Any Witness]
Let $(X,d)$ be a metric space, let \(x\in X\), and let \(a\in A\). Then
\[
  d(x,A)\leq d(x,a).
\]
\end{theorem*}
\hyperref[prf:fv-msml-distance-to-set-le-distance-to-point-of-mem]{\textit{Go to proof.}}
\end{remark*}

\begin{remark*}[Proof target]
\phantomsection
\label{thm:fv-msml-distance-to-set-eq-zero-of-mem}
\begin{theorem*}[Distance to a Set Vanishes at a Point of the Set]
Let $(X,d)$ be a metric space. If \(x\in A\), then
\[
  d(x,A)=0.
\]
\end{theorem*}
\hyperref[prf:fv-msml-distance-to-set-eq-zero-of-mem]{\textit{Go to proof.}}
\end{remark*}

\begin{remark*}[Formalization plan]
The first pass through these definitions is explicit about where \texttt{sInf}
and \texttt{sSup} enter the proof state. These topics are useful before balls
because open and closed balls quickly lead to boundedness, containment, and
estimates involving all points of a set. The current Lean names are
\texttt{distanceSet}, \texttt{distanceToSet}, \texttt{diameterSet}, and
\texttt{diameter}.
\end{remark*}

\subsection{Open and Closed Balls}

\begin{remark*}[Mathematical reminder]
\begin{definition*}[Open Ball]
Let $(X,d)$ be a metric space, let $x\in X$, and let $r>0$. The
\textbf{open ball} with center $x$ and radius $r$ is
\[
  B(x,r)=\{y\in X\mid d(x,y)<r\}.
\]
\end{definition*}
\end{remark*}

\begin{remark*}[Mathematical reminder]
\begin{definition*}[Closed Ball]
Let $(X,d)$ be a metric space, let $x\in X$, and let $r\geq 0$. The
\textbf{closed ball} with center $x$ and radius $r$ is
\[
  \overline{B}(x,r)=\{y\in X\mid d(x,y)\leq r\}.
\]
\end{definition*}
\end{remark*}

\begin{remark*}[Lean names]
In mathlib, the open ball is written
\[
  \texttt{Metric.ball x r}
\]
and the closed ball is written
\[
  \texttt{Metric.closedBall x r}.
\]
The active file \texttt{MetricBalls.lean} treats both as ordinary sets.
Membership is rewritten with \texttt{Metric.mem\_ball'} and
\texttt{Metric.mem\_closedBall'}, whose orientation matches the usual
\(d(x,y)\) notation used in hand proofs.
\end{remark*}

\begin{remark*}[Proof targets]
\phantomsection
\label{thm:fv-msml-center-mem-ball}
\begin{theorem*}[The Center Belongs to Every Positive Open Ball]
Let $(X,d)$ be a metric space. If $x\in X$ and $0<r$, then
\[
  x\in B(x,r).
\]
\end{theorem*}
\hyperref[prf:fv-msml-center-mem-ball]{\textit{Go to proof.}}
\[
  \texttt{center\_mem\_ball}
\]
is fully proved in Lean by applying \texttt{Metric.mem\_ball\_self}.
\end{remark*}

\begin{remark*}[Proof target]
\phantomsection
\label{thm:fv-msml-ball-subset-ball}
\begin{theorem*}[Open Balls Are Monotone in the Radius]
Let $(X,d)$ be a metric space. If $0<r\leq s$, then
\[
  B(x,r)\subseteq B(x,s).
\]
\end{theorem*}
\hyperref[prf:fv-msml-ball-subset-ball]{\textit{Go to proof.}}
\[
  \texttt{ball\_subset\_ball}
\]
is fully proved in Lean by introducing an arbitrary point, rewriting ball
membership with \texttt{Metric.mem\_ball'}, and composing the strict inequality
with the radius inequality.
\end{remark*}

\begin{remark*}[Proof target]
\begin{theorem*}[Closed Balls Are Monotone in the Radius]
Let $(X,d)$ be a metric space. If $0\leq r\leq s$, then
\[
  \overline{B}(x,r)\subseteq \overline{B}(x,s).
\]
\end{theorem*}
\end{remark*}

\begin{remark*}[Proof target]
\phantomsection
\label{thm:fv-msml-ball-subset-closed-ball}
\begin{theorem*}[Open Balls Sit Inside Closed Balls]
Let $(X,d)$ be a metric space. If $r\geq 0$, then
\[
  B(x,r)\subseteq \overline{B}(x,r).
\]
\end{theorem*}
\hyperref[prf:fv-msml-ball-subset-closed-ball]{\textit{Go to proof.}}
\[
  \texttt{ball\_subset\_closedBall}
\]
is fully proved in Lean by rewriting open-ball and closed-ball membership and
using \texttt{le\_of\_lt}.
\end{remark*}

\begin{remark*}[Proof target]
\phantomsection
\label{thm:fv-msml-ball-subset-ball-of-mem}
\begin{theorem*}[A Point Inside an Open Ball Has a Smaller Ball Inside It]
Let $(X,d)$ be a metric space. If $y\in B(x,r)$, then there exists
$\epsilon>0$ such that
\[
  B(y,\epsilon)\subseteq B(x,r).
\]
The usual choice is \(\epsilon=r-d(x,y)\).
\end{theorem*}
\hyperref[prf:fv-msml-ball-subset-ball-of-mem]{\textit{Go to proof.}}
\[
  \texttt{ball\_subset\_ball\_of\_mem}
\]
is fully proved in Lean with the choice \(\epsilon=r-d(x,y)\),
\texttt{dist\_triangle}, \texttt{linarith}, and \texttt{ring}.
\end{remark*}

\begin{remark*}[Formalization plan]
Each theorem above should first be stated against mathlib's existing
\texttt{Metric.ball} and \texttt{Metric.closedBall}. The early proofs should
use membership unfolding and elementary inequalities. The final theorem is the
first substantial metric-space proof in this track: it uses the triangle
inequality to show that any point close enough to \(y\) is still within radius
\(r\) of \(x\).
\end{remark*}

\subsection{Carrier, Typeclass, and Scratch Model}

\begin{remark*}[Mathematical carrier versus Lean type]
In the mathematical presentation, the symbol $X$ is the underlying set of
points. In Lean, this role is usually played by a type, for example
\texttt{X : Type u}. A metric-space assumption is then not a second component
manually paired with \texttt{X}; it is typeclass data attached to the type.
The local context has the shape
\[
  \texttt{variable \{X : Type u\} [MetricSpace X]}.
\]
This says that \texttt{X} is an arbitrary type equipped with a mathlib
metric-space structure.
\end{remark*}

\begin{remark*}[Why not define the main metric ourselves]
One can understand the mathematics by imagining a structure with a distance
function and fields proving the metric axioms. That is a useful mental model,
but it is not the implementation strategy for the mathlib track. Defining a
local metric structure as the main interface would create a parallel API and
force every theorem to be translated back to Mathlib. For this track, theorem
files use Mathlib's existing \texttt{MetricSpace} class, its \texttt{dist}
notation, and the theorem names already proved for that class.
\end{remark*}

\begin{remark*}[Scratch modeling of the space and metric]
The handwritten model remains important as scaffolding for learning. A direct
Lean model separates the carrier from the distance function. If
\texttt{Point} is the Lean type playing the role of the mathematical set \(X\),
then a metric can be modeled as data of the following shape:
\begin{verbatim}
structure ScratchMetric (Point : Type u) where
  distance : Point -> Point -> Real
  distance_nonnegative :
    forall x y : Point, 0 <= distance x y
  distance_eq_zero_iff_equal :
    forall x y : Point, distance x y = 0 <-> x = y
  distance_symmetric :
    forall x y : Point, distance x y = distance y x
  distance_triangle :
    forall x y z : Point,
      distance x z <= distance x y + distance y z
\end{verbatim}
The field \texttt{distance} is the proposed distance function. The remaining
fields are proof obligations. Constructing a value of type
\texttt{ScratchMetric Point} therefore means supplying both a formula and
proofs that the formula satisfies the metric laws.
\end{remark*}

\begin{remark*}[Scratch metric-space wrapper]
One may then package the metric as the chosen metric on a carrier:
\begin{verbatim}
structure ScratchMetricSpace (Point : Type u) where
  metric : ScratchMetric Point
\end{verbatim}
This wrapper expresses an important modeling point: the same underlying type
can carry different metrics. Changing the metric changes the structure while
leaving the carrier type fixed.

In the Mathlib track, this wrapper is not the main interface. It is a
controlled learning model for understanding part of what Mathlib's
\texttt{MetricSpace} typeclass supplies. The bridge
\texttt{ScratchMetricSpace.toMathlibMetricSpace} also shows why the conversion
is not automatic: Mathlib's metric-space hierarchy carries compatible
topological and uniform structure, not just a distance function.
\end{remark*}

\begin{remark*}[Local Lean context]
For theorem statements about arbitrary metric spaces in the main development,
the local context has the shape:
\begin{verbatim}
variable {X : Type u} [MetricSpace X]
\end{verbatim}
There is no explicit bundled metric parameter. Lean finds the metric-space
instance by typeclass search, and \texttt{dist x y} uses that instance.
\end{remark*}

\begin{example*}[Repackaging a scratch metric axiom]
In the scratch model, a theorem such as
\begin{verbatim}
theorem equal_of_distance_eq_zero
    {x y : Point} (hypothesis : M.distance x y = 0) :
    x = y :=
  (M.metric.distance_eq_zero_iff_equal x y).mp hypothesis
\end{verbatim}
does not add a new axiom. It repackages one direction of an axiom field already
stored in \texttt{M.metric}. If
\[
  \texttt{h : A <-> B},
\]
then \texttt{h.mp} has type \texttt{A -> B}, and \texttt{h.mpr} has type
\texttt{B -> A}. This is the basic Lean technique behind extracting one
direction from an if-and-only-if proof.
\end{example*}

\begin{remark*}[Axioms versus Lean axioms]
The metric requirements are called axioms because they are defining laws of the
mathematical structure. In Lean, they should not be introduced as global
logical axioms. In the scratch structure, they are fields whose values must be
supplied when constructing an example. In the Mathlib-first files, these facts
are supplied by the existing \texttt{MetricSpace} instance and used through
API lemmas such as \texttt{dist\_triangle}.
\end{remark*}

\begin{remark*}[Namespace]
The corresponding Lean files for this track live under the namespace
\[
  \texttt{LRA.VolumeVII.WithMathlib}.
\]
The initial module for metric spaces is
\[
  \texttt{LRA.VolumeVII.WithMathlib.MetricSpaces}.
\]
\end{remark*}

\subsection{Real-Line Metric}

\begin{example*}[Absolute-value metric as a proof target]
The real line $\mathbb{R}$ carries the usual absolute-value metric
\[
  d(x,y)=|x-y|.
\]
If this metric were being constructed by hand, the proof obligations would be
nonnegativity, identity of indiscernibles, symmetry, and the triangle
inequality:
\[
  0\le |x-y|,\qquad
  |x-y|=0\iff x=y,\qquad
  |x-y|=|y-x|,
\]
and
\[
  |x-z|\le |x-y|+|y-z|.
\]
In the mathlib track, these obligations are not new axioms. They are either
already part of the existing instance for \(\mathbb{R}\) or available as
library theorems about absolute value.
\end{example*}

\begin{example*}[Beginning the scratch construction]
For learning purposes, the scratch construction begins by fixing the distance
formula and exposing the four proof obligations:
\begin{verbatim}
/-- The usual absolute-value scratch metric on the real line. -/
def realScratchMetric : ScratchMetric Real where
  distance x y := |x - y|
  distance_nonnegative := by
    sorry
  distance_eq_zero_iff_equal := by
    sorry
  distance_symmetric := by
    sorry
  distance_triangle := by
    sorry
\end{verbatim}
At this stage the formula is fixed. The remaining work is to prove that this
formula satisfies each metric axiom.
\end{example*}

\begin{remark*}[Theorems that explain the proof]
The relevant library facts for the hand proof include:
\begin{itemize}
  \item \texttt{abs\_nonneg}, expressing $0\le |a|$;
  \item \texttt{abs\_eq\_zero}, expressing that an absolute value is zero
        exactly when its argument is zero;
  \item \texttt{abs\_sub\_comm}, expressing $|x-y|=|y-x|$;
  \item \texttt{abs\_add\_le}, expressing $|a+b|\le |a|+|b|$.
\end{itemize}
The triangle inequality also needs the algebraic bridge
\[
  x-z=(x-y)+(y-z).
\]
This is the common Lean pattern: a theorem may contain the right mathematics,
but the goal must first be put into the syntactic form where that theorem
applies.
\end{remark*}

\begin{example*}[Nonnegativity]
The nonnegativity field asks for a proof of \(0\le |x-y|\) for arbitrary real
numbers \(x\) and \(y\). Mathlib's theorem \texttt{abs\_nonneg} proves
\(0\le |a|\) for any real number \(a\). The Lean proof is:
\begin{verbatim}
distance_nonnegative := by
  intro x y
  exact abs_nonneg (x - y)
\end{verbatim}
After \texttt{intro x y}, the goal is exactly
\texttt{0 <= |x - y|}. Applying \texttt{abs\_nonneg} to the real number
\texttt{x - y} closes the goal.
\end{example*}

\begin{example*}[Identity of indiscernibles]
The identity field asks for a proof of
\[
  |x-y|=0\iff x=y.
\]
Because the goal is an if-and-only-if statement, \texttt{constructor} splits
the proof into the forward direction and the reverse direction:
\begin{verbatim}
distance_eq_zero_iff_equal := by
  intro x y
  constructor
  -- =>
  . intro hypothesis
    have difference_eq_zero : x - y = 0 :=
      abs_eq_zero.mp hypothesis
    linarith
  -- <=
  . intro hypothesis
    rw [hypothesis, sub_self]
    exact abs_zero
\end{verbatim}
The forward branch uses \texttt{abs\_eq\_zero.mp} to turn
\texttt{|x - y| = 0} into \texttt{x - y = 0}. The arithmetic step
\texttt{linarith} then proves \texttt{x = y}. The reverse branch rewrites
\texttt{x} to \texttt{y}, simplifies \texttt{y - y} by \texttt{sub\_self},
and finishes with \texttt{abs\_zero}.
\end{example*}

\begin{example*}[Symmetry]
The symmetry field asks for \( |x-y|=|y-x| \). Mathlib names this fact
\texttt{abs\_sub\_comm}:
\begin{verbatim}
distance_symmetric := by
  intro x y
  exact abs_sub_comm x y
\end{verbatim}
This is an example where the mathematical statement and the library theorem
match directly after the variables are introduced.
\end{example*}

\begin{example*}[Triangle inequality calculation]
The hand proof of the real-line triangle inequality has the Lean shape
\begin{verbatim}
calc
  |x - z| = |(x - y) + (y - z)| := by
    ring_nf
  _ <= |x - y| + |y - z| := by
    exact abs_add_le (x - y) (y - z)
\end{verbatim}
The first line is algebraic normalization. The second line applies the
absolute-value triangle inequality to the two terms \(x-y\) and \(y-z\).
\end{example*}

\begin{example*}[Complete scratch construction]
Putting the four fields together gives the completed scratch construction:
\begin{verbatim}
/-- The usual absolute-value scratch metric on the real line. -/
def realScratchMetric : ScratchMetric Real where
  distance x y := |x - y|

  distance_nonnegative := by
    intro x y
    exact abs_nonneg (x - y)

  distance_eq_zero_iff_equal := by
    intro x y
    constructor
    -- =>
    . intro hypothesis
      have difference_eq_zero : x - y = 0 :=
        abs_eq_zero.mp hypothesis
      linarith
    -- <=
    . intro hypothesis
      rw [hypothesis, sub_self]
      exact abs_zero

  distance_symmetric := by
    intro x y
    exact abs_sub_comm x y

  distance_triangle := by
    intro x y z
    calc
      |x - z| = |(x - y) + (y - z)| := by
        ring_nf
      _ <= |x - y| + |y - z| := by
        exact abs_add_le (x - y) (y - z)
\end{verbatim}
The purpose of this construction is pedagogical. The mathlib track should then
translate this understanding into use of mathlib's existing
\texttt{MetricSpace} instance for \(\mathbb{R}\).
\end{example*}

\subsection{Discrete Metric}

\begin{example*}[Discrete metric]
Every type with decidable equality carries the discrete metric
\[
  d(x,y)=
  \begin{cases}
    0, & x=y,\\
    1, & x\ne y.
  \end{cases}
\]
In Lean, the definition uses an \texttt{if}-expression and therefore requires
\[
  \texttt{[DecidableEq X]}.
\]
Proofs about this metric usually proceed by case splitting with
\texttt{by\_cases} and simplifying the relevant \texttt{if}-expressions with
\texttt{simp}. The scratch definitions are named
\texttt{discreteScratchMetric} and \texttt{discreteScratchMetricSpace}.
\end{example*}

\subsection{Euclidean Plane}

\begin{example*}[Mathlib Euclidean space]
The Mathlib-first object for the real Euclidean plane is
\[
  \texttt{EuclideanSpace Real (Fin 2)}.
\]
The Lean file abbreviates this as
\begin{verbatim}
abbrev cartesianSpace : Type :=
  EuclideanSpace Real (Fin 2)
\end{verbatim}
Mathlib supplies the metric-space instance. The chapter does not build a new
main metric-space structure for \(\mathbb{R}^2\).
\end{example*}

\begin{example*}[A concrete Euclidean distance calculation]
The file defines the origin and the point \((1,1)\) as Euclidean-space points:
\begin{verbatim}
noncomputable def cartesianOrigin : cartesianSpace :=
  (EuclideanSpace.equiv (Fin 2) Real).symm ![0, 0]

noncomputable def cartesianOneOne : cartesianSpace :=
  (EuclideanSpace.equiv (Fin 2) Real).symm ![1, 1]
\end{verbatim}
Then it proves
\begin{verbatim}
lemma dist_cartesianOrigin_cartesianOneOne :
    dist cartesianOrigin cartesianOneOne = Real.sqrt 2 := by
  rw [EuclideanSpace.dist_eq, Fin.sum_univ_two]
  simp [cartesianOrigin, cartesianOneOne, EuclideanSpace.equiv, Real.dist_eq]
  ring_nf
\end{verbatim}
The proof rewrites Mathlib's Euclidean distance with
\texttt{EuclideanSpace.dist\_eq}, expands the finite sum over
\texttt{Fin 2}, simplifies the coordinates, and normalizes the arithmetic.
\end{example*}

\begin{remark*}[Using mathlib's Euclidean space]
For the triangle inequality, the most maintainable mathlib-based route is to
connect the coordinate formula to mathlib's Euclidean-space distance and then
use the existing theorem \texttt{dist\_triangle}. This is not a direct proof
from square-root algebra; it is an interoperation proof. The value of the
exercise is learning how to bridge a local coordinate expression to a library
object that already has the required theorem.
\end{remark*}

\begin{example*}[Bridge to mathlib Euclidean space]
The relevant import for the bridge is:
\begin{verbatim}
import Mathlib.Analysis.InnerProductSpace.PiL2
\end{verbatim}
Mathlib represents the real Euclidean plane as
\[
  \texttt{EuclideanSpace Real (Fin 2)}.
\]
Points represented as pairs can be converted into this space by sending
\texttt{p : Real \(\times\) Real} to the length-two coordinate vector
\texttt{![p.1, p.2]}:
\begin{verbatim}
noncomputable def toE (p : Prod Real Real) :
    EuclideanSpace Real (Fin 2) :=
  (EuclideanSpace.equiv (Fin 2) Real).symm ![p.1, p.2]
\end{verbatim}
The bridge lemma states that the coordinate square-root formula is mathlib's
distance after this conversion:
\begin{verbatim}
lemma distance_eq_euclidean (p q : Prod Real Real) :
    Real.sqrt ((p.1 - q.1) ^ 2 + (p.2 - q.2) ^ 2)
      = dist (toE p) (toE q) := by
  rw [EuclideanSpace.dist_eq, Fin.sum_univ_two]
  simp [toE, EuclideanSpace.equiv, Real.dist_eq, sq_abs]
\end{verbatim}
Once this bridge is available, the triangle inequality is reduced to
\texttt{dist\_triangle}:
\begin{verbatim}
distance_triangle := by
  intro x y z
  rw [distance_eq_euclidean, distance_eq_euclidean,
    distance_eq_euclidean]
  exact dist_triangle (toE x) (toE y) (toE z)
\end{verbatim}
This is the second full pass through proving a metric: first by direct
absolute-value calculations on \(\mathbb{R}\), and then by a combination of
coordinate algebra and library interoperation on \(\mathbb{R}^2\).
\end{example*}
